\documentclass[11pt]{article}

\usepackage[margin=1in]{geometry}
\usepackage{amsmath,amssymb}
\usepackage{enumitem}
\usepackage[colorlinks=true,linkcolor=blue,urlcolor=blue,citecolor=blue]{hyperref}
\emergencystretch=2em

\newcommand{\HC}{\operatorname{HC}}

\title{A Literature-Implied Negative Answer to Question 4\\
on Hypercyclic-Algebra Generators}
\author{Math discovery status packet}
\date{August 17, 2026}

\begin{document}
\maketitle

\section{Status and source question}

This is a full \emph{literature-implied negative answer}, not a new result.
Abbar--Costa Jr., \emph{Hypercyclic algebras for weighted shifts on trees},
arXiv:2411.14609v1, asks in Question 4 (printed page 31): if a continuous
linear map \(T:X\to X\) on an \(F\)-space supports a hypercyclic algebra, must
the set of elements generating a hypercyclic algebra be residual, or at least
dense?

Bayart--Costa Jr.--Papathanasiou, \emph{Baire theorem and hypercyclic
algebras}, arXiv:1910.05000, Remark 6.7 (printed pages 50--51), had already
given an example in which the generator set is nonempty and nowhere dense.
Thus both questions have negative answers.

\section{Exact identification}

Let \(T\) be a hypercyclic operator on a Banach space \(X\), choose
\(h\in\HC(T)\), and choose a nonzero \(y\notin\HC(T)\) linearly independent
of \(h\).  Take \(f\in X^*\) with
\(f(h)=0\) and \(f(y)=1\), and define
\[
        z\cdot w=f(z)f(w)y \qquad(z,w\in X).
\]
This is a continuous commutative associative multiplication; associativity
uses \(f(y)=1\).  After an equivalent rescaling of the norm if necessary, it
is a Banach-algebra norm.

Because \(h^2=0\), the non-unital algebra generated by \(h\) is
\(\operatorname{span}\{h\}\), so every nonzero element of it is hypercyclic.
On the other hand, if \(f(z)\ne0\), then
\(z^2=f(z)^2y\) is a nonzero non-hypercyclic vector.  Consequently Remark 6.7
identifies the generator set exactly as
\[
 \bigl\{z\in X:\mathcal A(z)\setminus\{0\}\subset\HC(T)\bigr\}
       =\HC(T)\cap\ker f.
\]
It is nonempty because it contains \(h\), and it is nowhere dense because it
lies in the proper closed hyperplane \(\ker f\).

For a concrete underlying pair one can take the Rolewicz operator \(2B\) on
\(\ell^2\): it is hypercyclic and has nonzero finitely supported
non-hypercyclic vectors.  Hence the construction gives a literal countermodel
to the universal assertion in Question 4.

\section{Provenance and scope}

The source paper cites the 2021 article in its bibliography, but its Question
4 does not mention Remark 6.7 or explain the conflict with the question's
general wording.  The implication above is therefore recorded rather than
claimed as novel.

The example engineers the multiplication.  It does \emph{not} decide whether
generators are dense or residual for the fixed coordinatewise product on the
tree sequence algebras studied in arXiv:2411.14609, nor does it address that
paper's Questions 1--3 or 5.

\section{References}

\begin{itemize}[leftmargin=*]
\item A. Abbar and F. Costa Jr., \emph{Hypercyclic algebras for weighted
shifts on trees}, arXiv:2411.14609v1 (2024), Question 4, p.~31.
\item F. Bayart, F. Costa Jr., and D. Papathanasiou, \emph{Baire theorem and
hypercyclic algebras}, Advances in Mathematics 376 (2021), 107419;
arXiv:1910.05000, Remark 6.7, pp.~50--51.
\end{itemize}

\end{document}
